\documentclass[a4paper,12pt]{article}
\usepackage[T1]{fontenc}
\usepackage[utf8]{inputenc}  % à garder si tu compiles avec pdflatex
\usepackage{lmodern}
\usepackage{graphicx}
\usepackage{listings}
\usepackage{xcolor}
\usepackage{hyperref}
\usepackage{geometry}
\usepackage{abstract}
\usepackage{titling}
\geometry{margin=2.5cm}

\lstset{
    language=C++,
    basicstyle=\ttfamily\footnotesize,
    keywordstyle=\color{blue},
    commentstyle=\color{gray},
    stringstyle=\color{orange},
    showstringspaces=false,
    frame=single,
    numbers=left,
    numberstyle=\tiny,
    numbersep=5pt,
    breaklines=true,
    breakatwhitespace=true,
    columns=fullflexible,
    inputencoding=utf8,
    literate=
      {é}{{\'e}}1 {è}{{\`e}}1 {ê}{{\^e}}1 {ë}{{\"e}}1
      {à}{{\`a}}1 {â}{{\^a}}1 {ä}{{\"a}}1
      {ù}{{\`u}}1 {û}{{\^u}}1 {ü}{{\"u}}1
      {ô}{{\^o}}1 {ö}{{\"o}}1 {î}{{\^i}}1 {ï}{{\"i}}1
      {ç}{{\c{c}}}1
      {É}{{\'E}}1 {È}{{\`E}}1 {Ê}{{\^E}}1 {Ë}{{\"E}}1
}

\begin{document}

\begin{titlepage}
    \centering
    \vspace*{\fill}

    {\Huge \bfseries Compte Rendu de Projet \\[0.5em]
    Simulation de Bateau en C++}

    \vspace{2cm}
    {\Large RULLAC Éloïse, FAUCON Léo \par}
    {\large Université de Toulon \par}
    {\large UE622 - Programmation Orientée Objet \par}
    \vspace{1.5cm}
    {\large \today}

    \vspace*{\fill}
\end{titlepage}

\section{Introduction}
Le projet consistait à simuler une course de bateau via la bibliothèque SFML et libphysics (bibliothèque locale au module UE622). Ce travail s'inscrit dans le cadre de l'UE622 à l'Université de Toulon. L'objectif pédagogique principal est l'application des notions de POO en C++.

\section {Fonctionnalitées}
Avec ses menus, ses bateaux aussi customisables que notre concurrent officiel Sea of Thieves et
son réalisme à couper le souffle, notre projet va changer votre vision du jeu vidéo maritime.
Notre projet est donc : 
\begin{itemize}
    \item Jouable de 0 (oui 0) à 4.
    \item Doté de bot.
    \item Pas réaliste, car les calculs sont fait par nous-même,
     sans l'aide de la bibliothèque libphysics.
    Par exemple il est possible de faire un dash (c'est volontaire)
    via la pression simultanée de la touche "avancer" et "reculer",
    l'accélération du bateau va s'accumuler,
    et lors du relâchement de la touche "reculer", le bateau profiterai alors du boost accumulé.
\end{itemize}

Il possède :
\begin{itemize}
    \item \texttt{Un écran de début} où est possible:
    \begin{itemize}
        \item L'ajout de bateaux (bot ou joueur).
        \item La configuration de ces bateaux (taille, couleur).
    \end{itemize}

    \item \texttt{Un écran de fin} où le choix de recommencer est permis (via un bouton ou la touche espace).
    \item \texttt{La possibilité de tirer des boulets de canon}:
    si beau que cela puisse paraître la hitbox est honteusement gérée.
    La bibliothèque SFML offre la possibilité de vérifier si deux sf::Shape
    s'intersectent via la méthode :
    \[
    \texttt{sf::FloatRect::intersects(sf::FloatRect)}
    \]
    Or comme cela semble être intuitif dans le nom,
    cette méthode gère uniquement les intersections des rectangles.
    Et par nous ne savons quelle magie,
    il est possible de récupérer un \texttt{sf::FloatRect} via
    \texttt{sf::Shape::getGlobalBounds()} donc avec des formes complexes du genre 
    \texttt{CircleShape}, \texttt{ConvexShape}, ..., mais est terrible. C'est ainsi pourquoi nous avons dû
    faire tirer les boulets bien plus loin pour éviter que les bateaux n'explosent de leur propres tirs.
\end{itemize}

\section{Structuration du code}
\subsection{POO}
Le projet est composé de plusieurs classes :
\begin{itemize}
    \item \textbf{Boutons} :
    \begin{itemize}
        \item \texttt{BareButton} : bouton sans texte, juste un rectangle
        \item \texttt{Button} : bouton classique avec texte, hérite de \texttt{BareButton}
        \item \texttt{RadioButtonG\textless int nb\_button, typename ButtonT, typename ValueT\textgreater} : système de boutons radio configurable par template
        \item \texttt{RadioButton\textless int nb\_button, typename ValueT\textgreater} : bouton radio classique, hérite de \texttt{RadioButtonG\textless int nb\_button, Button, ValueT\textgreater}
    \end{itemize}
    
    \item \textbf{Éléments de jeu} :
    \begin{itemize}
        \item \texttt{Entity\textless typename T\textgreater} : Concept global d'une entité, hérite de \texttt{sf::Drawable}
        \item \texttt{Obstacle} : obstacle fixe sur le parcours, hérite de \texttt{Entity\textless sf::RectangleShape\textgreater}
        \item \texttt{Cannonball} : projectile tiré par les bateaux, hérite de \texttt{Entity\textless sf::CircleShape\textgreater}
        \item \texttt{FinishLine} : ligne d'arrivée, hérite de \texttt{Entity\textless sf::RectangleShape\textgreater}
        \item \texttt{Screamer} : effet visuel et sonore lancé lors de la destruction d'un bateau, hérite de \texttt{sf::Drawable}
    \end{itemize}
    
    \item \textbf{Bateaux} :
    \begin{itemize}
        \item \texttt{Boat} : bateau générique, hérite de \texttt{Entity\textless sf::ConvexShape\textgreater}
        \item \texttt{BoatAI} : bateau contrôlé par une IA, hérite de \texttt{Boat}
        \item \texttt{BoatPlayer} : bateau contrôlé par un joueur, hérite de \texttt{Boat}
    \end{itemize}
    
    \item \textbf{Interface} :
    \begin{itemize}
        \item \texttt{StartScreen} : écran d'accueil du jeu, hérite de \texttt{sf::Drawable}
        \item \texttt{EndScreen} : écran de fin de partie, hérite de \texttt{sf::Drawable}
        \item \texttt{BoatOption} : options de configuration d'un bateau, hérite de \texttt{sf::Drawable}
        \item \texttt{Players} : gestion de tous les joueurs (humains ou IA)
        \item \texttt{Game} : classe principale gérant la boucle du jeu
    \end{itemize}
\end{itemize}

\subsection{Architecture du Projet}
Pour ce projet, nous avons décidé que toutes les entitées ne seraient constituées que d'une forme, ou \texttt{sf::Shape}, ainsi, la détection de collision via SFML est aisement vérifiée via le code suivant dans la classe \texttt{Entity} :

\begin{lstlisting}
// T doit hériter de sf::Shape
template <typename T>
class Entity: public sf::Drawable {
public:
    template <typename U>
    inline bool hasCrossed(const Entity<U>& e) const
    {
        return shape.getGlobalBounds().intersects(e.get_shape().getGlobalBounds());
    }
    (...)
protected:
    T shape;
};
\end{lstlisting}
Le code est évidemment sectionné, symbolisé par le "(...)", où sont les méthodes basiques, tirées de celles de sf::Shape, du genre setPosition, setSize, ...

\section {Les difficultées et solutions}
Nous avons rencontré plusieurs difficultées lors de la conception du projet.
\subsection{Le polymorphisme}
    Ce concept de POO a été central dans notre projet.
    Nous l'avons modularisé au maximum,
    jusqu'à créer le concept entier d'un radiobutton en classe pour
    deux pauvres instances.
\subsection{La gestion des instances temporaires de classes}
    Celle-ci est un peu plus compliqué à expliquer. Nous sommes passés par plusieurs méthodes.
\subsubsection{La méthode ObstacleVect}
Pour la première méthode, nous avons fait une liste basique de boulets, via la classe
\texttt{ObstacleVect} qui hérite de \texttt{std::vect\textless Obstacle\textgreater}
et gère les instances des obstacles avec les méthodes \texttt{inst} et \texttt{reset}. 
\subsubsection{La méthode Cannonball}
Pour la deuxième méthodes, nous voulions tenter un système plus performant. Donc nous avons fait ceci :
    \begin{lstlisting}
class Cannonball: public Entity<sf::CircleShape> {
private:
    Cannonball(const sf::Vector2f& starting_pos, const sf::Vector2f& pdirection);

    (...)

public:
    static std::vector<std::shared_ptr<Cannonball>> all;
    Cannonball() = delete;
    Cannonball(const Cannonball&) = delete;
    Cannonball& operator=(const Cannonball&) = delete;

    static Cannonball& create(const sf::Vector2f& starting_pos, const sf::Vector2f& pdirection);

    static void drawAll(sf::RenderTarget& target);
    inline static void delInst(const std::vector<std::shared_ptr<Cannonball>>::iterator& it);

    (...)
};
    \end{lstlisting}
    Le constructeur, la copie et l'affection sont supprimés au profit
    de la méthode de classe create, qui alloue un boulet, insère son smart pointer dans l'attribut
    de classe all, et retourne l'instance crée.
    Cela permet de gérer toutes les instances via les méthodes de classes :
    \texttt{drawAll} et \texttt{delInst}
\subsubsection{La méthode BareButton}
Cette dernière méthode consiste en ce code:
\begin{lstlisting}
class BareButton: public sf::Drawable {
public:
    static std::vector<BareButton*> all;
    static void checkAllClicks(const sf::Vector2f& mousePos);

    BareButton();
    virtual ~BareButton() { all.erase(std::find(all.begin(), all.end(), this)); }
    BareButton(const sf::Vector2f& pos);
    BareButton(const BareButton&) = delete;
    BareButton& operator=(const BareButton&) = delete;
};
\end{lstlisting}
Comme on peut le remarquer, BareButton possède un attribut de classe all
qui stock tous les pointeurs des instances. Les pointeurs sont ajoutés à cet attribut dans les constructeurs.
La copie et l'affection sont supprimées par précaution,
pour éviter de se faire
surprendre par d'éventuels soucis de pointeurs.

\section {La répartition des tâches}
Ce fut comme ceci :
\begin{itemize}
    \item \texttt{Éloïse} s'est occupée du fondement du projet c'est à dire
    le fonctionnement des bateaux (bot et joueurs), obstacles, boulets de canons.
    \item \texttt{Léo} s'est occupé de la structuration du projet (polymorphisme, gestion des instances, modularité)
    et de l'interface (StartScreen, EndScreen, BareButton...).
\end{itemize}

\section {Retour d'experience sur le travail en groupe}
La répartition des tâches pour ce projet fut complexe,
étant donné le manque de temps dû aux autres projets et partiels à venir,
il était compliqué de réaliser des idées poussées pouvant être
codées séparément. Nous étions souvent ensemble lors de l'avancement du projet.

\vfill
\noindent\textbf{Références} :
\begin{itemize}
    \item Documentation C++ : \url{https://en.cppreference.com/}
    \item Cours UE622 - Université de Toulon
\end{itemize}

\end{document}
